{ % Concurrency Theatre companion, audited 24 October 2026
  \section{Concurrency Theatre: an interactive companion}
  \label{concurrencyTheatre:section}

  \subsection{Provenance and purpose}

  \par This companion connects selected ideas in the original teaching material to reproducible browser experiments. The archival revision is dated 19 February 2020, commit \texttt{1c53d6eb08a7d0041a6f54728b6ba90a59507454} in the repository \texttt{SkyProse/ParallelProgrammingTex}. The original attribution to V. V. Sosnin, P. V. Balakshin, D. S. Shilko and ITMO University is preserved. This addition, audited on 24 October 2026, does not claim to be a new official ITMO edition. Original C++ listing files are retained unchanged; a teaching adaptation is identified separately from its source.

  \par The learning cycle is: predict an outcome, execute a stated model or experiment, explain the causal difference, and solve a new transfer question. A successful animation is not evidence that a real C++ program is correct. A successful test on one schedule is not evidence that every schedule succeeds.

  \subsection{Three kinds of evidence}

  \begin{enumerate}
    \item \textbf{Model.} A deterministic state machine applies explicit rules. Its events have a logical order, not measured durations. A cache-line transfer counter is a count within that model, not a hardware performance-counter reading.
    \item \textbf{Measured.} A separate experiment actually executes in the current environment. It records the input, correctness checks, timing boundary, repetitions and relevant capabilities. Opening an experiment page does not create a measurement.
    \item \textbf{Imported.} A saved browser run or native program result is loaded for inspection. Structural validation does not authenticate its origin or repeat its measurements on the current machine. The imported origin and timing boundary remain visible.
  \end{enumerate}

  \subsection{Six scenes and their boundaries}

  \paragraph{A fork in time: deadlock.}
  \par The adapted example uses two threads, two ordinary locks and one iteration. Each thread executes five instructions: acquire its first lock, acquire its second lock, perform work, release the second lock and release the first. A blocked acquisition does not advance the program counter. If T0 acquires A and T1 acquires B while their requested orders differ, both next acquisitions are disabled. The wait cycle is sufficient to identify a deadlock in this finite model; no timeout is required.
  \par Rewinding and choosing a different enabled thread can produce a successful execution of the same program. Changing both programs to acquire A before B removes the wait cycle in this model. These are different interventions. A graph of reachable states is an exhaustive result only for the stated finite state machine, not for arbitrary C++, an operating system scheduler or additional external resources. Source: \path{listings/deadlockExample.cpp}.
  \par\textit{Transfer question.} A thread waits for a lock for 10 ms, but the owner can still run and release it. Explain why the elapsed time alone does not establish a deadlock.

  \paragraph{The contract of a race.}
  \par The original examples contrast reading a changing variable with printing a snapshot named \texttt{old\_a}. The browser scene instead declares a sequentially consistent toy machine with two read--increment--write operations on a value initially equal to zero. Its unsafe mode can lose an update. Its atomic mode preserves the contract of two shared increments. Its snapshot mode computes two local results from the initial value without changing the shared counter.
  \par This adaptation must not be confused with an enumeration of C++ data-race outcomes. A data race on an ordinary C++ object is undefined behavior; a visual list of a few plausible values does not constrain it. A snapshot can be the correct solution to an initial-state reporting task, but it changes the contract of a shared-counter task. Sources: \path{listings/raceConditionalExample1.cpp} and \path{listings/raceConditionalExample2.cpp}; C++ rules: \url{https://eel.is/c++draft/intro.races}.
  \par\textit{Transfer question.} Which solution fits two independent calculations from the state before parallel work begins? Which fits two changes that must both appear in one shared counter?

  \paragraph{The memory microscope.}
  \par The scene assigns one byte each to fields A and B. AoS interleaves the two fields; SoA places the arrays separately with a line-aligned start for each array; padding gives each scalar an aligned slot of at least one configured line. The requested gap and resulting memory footprint are explicit. Line size is a model parameter and must not be presented as a universal processor constant.
  \par The simplified coherence model distinguishes remote copies affected by access to the same scalar from copies affected by a different scalar in the same line. Read-only sharing does not count as false sharing. The model does not implement full MESI, cache capacity, prefetching, NUMA, a compiler or hardware timestamps. A native packed-versus-padded experiment is a separate measurement with its own assumptions. Source: \path{listings/falseSharingExample.cpp}; practical context: \url{https://www.kernel.org/doc/html/latest/kernel-hacking/false-sharing.html}.
  \par\textit{Transfer question.} Why does placing two different fields on separate lines not remove the need to synchronize updates of the same field?

  \paragraph{The CAS relay.}
  \par The queue scene admits two producers and unique new nodes, with append operations only. A successful strong CAS that changes a node's \texttt{next} link from null to the new node is the insertion's linearization point in this model. Advancing a lagging \texttt{tail} pointer is useful but can be completed by another thread. A failed CAS updates \texttt{expected} with the observed value.
  \par The optional spurious-failure control applies only to the weak CAS used while helping the tail. It is not applied to the strong CAS that links a new node. The scene does not implement dequeue, safe memory reclamation or an ABA defence, and it does not prove correctness for arbitrary C++ memory orders. Source: \path{listings/lockFreeQueueEnqueue.cpp}; CAS semantics: \url{https://eel.is/c++draft/atomics.types.generic}.
  \par\textit{Transfer question.} If the next-link CAS succeeds but the final tail CAS fails, what has already happened, and what can another thread help finish?

  \paragraph{The GPU loom.}
  \par The OpenCL chapter supplies the historical idea of a host launching indexed work on a device. The new scene uses WGSL and WebGPU, not the OpenCL API. A real shader returns global, local and workgroup identifiers. Subgroup values are used only when the necessary features are available. A requested workgroup size does not specify a physical SIMD width, and a subgroup size must not be assumed to equal 32 or 64.
  \par The resulting diagram is a representation of returned identifiers, not a hardware execution timeline. Missing features, an unavailable adapter and device loss are explicit absence of results. New WGSL capabilities and HTML-in-Canvas are not required to complete the basic lesson. Source: \path{section3_4.tex}; language reference: \url{https://www.w3.org/TR/2026/CRD-WGSL-20261021/}.
  \par\textit{Transfer question.} Cover 130 array elements using workgroups of 64 invocations. Explain the number of groups, bounds checks and, for an algorithm using barriers, the need to obey WGSL participation rules.

  \paragraph{The performance observatory.}
  \par This scene extends Laboratory Research No. 4. Browser JS, Web Workers and available WebGPU execute the same deterministic unsigned 32-bit transform and reduction. Every output value and the sum modulo $2^{32}$ are checked against a reference before speed is interpreted. The measurement loop pauses scene rendering; this reduces a source of interference but does not isolate the machine from other activity.
  \par Cold measurements include backend preparation and the first run. Warm measurements exclude preparation. The browser's end-to-end wall-clock interval includes the relevant output allocation, communication or upload, execution and readback; correctness validation is outside that interval. An optional GPU timestamp measures a different, device-side interval and is reported separately. A native OpenMP runner uses a separate process and a compute-only boundary with preallocated arrays. These numbers cannot be combined into a single speedup without aligning the protocol.
  \par A Student t interval estimates the sample mean under assumptions about the observations. The 95\% level refers to long-run coverage of the procedure, not a probability that a particular fixed interval contains the mean. Consecutive browser measurements may be dependent, so the interval is exploratory evidence rather than an automatic significance test. Preserve samples, sample size, the minimum, median, mean, sample standard deviation and interval method. Reference: \url{https://www.itl.nist.gov/div898/handbook/eda/section3/eda352.htm}.
  \par\textit{Transfer question.} Why is comparing a cold CPU duration with a warm GPU kernel-only duration insufficient to establish application-level acceleration?

  \subsection{Current standards and practical use}

  \par At the audit date, OpenMP 6.0 is final and OpenMP 6.1 TR15 is a public comment draft. A published specification does not imply full support in every compiler or on every device. Preserve exact compiler and runtime versions, flags and target information with a native experiment. The authoritative publication list is \url{https://www.openmp.org/specifications/}; the draft announcement is \url{https://www.openmp.org/home-news/public-comment-draft-6-1/}.

  \par WebGL 2 and WebGPU are separate APIs. A graphics fallback can preserve the scene while WebGPU compute is unavailable. Accessible text, keyboard controls and the two-dimensional state representation must remain useful regardless of a preferred renderer. Reduced motion changes presentation, not the semantics of an experiment.

  \par The web companion resides in \path{web/}. The teacher guide, source register, architecture and detailed errata reside in \path{docs/}. The project does not infer a student's understanding from watching an animation: each scene includes a challenge and a different transfer question. In a report, state the intended contract, provide a trace or measurement, explain the relevant cause, and identify the limit of the conclusion.
}
